\section{本章のまとめ}
\label{sec:quantum-subroutines-summary}

本章では，量子アルゴリズムを構成する重要なサブルーチンとして，
量子フーリエ変換，量子位相推定，振幅増幅を説明した．

量子フーリエ変換は，計算基底の振幅を離散フーリエ変換するユニタリ演算である．
この演算は，アダマールゲートと制御位相ゲートから構成できる．
量子位相推定では，制御$U^{2^r}$演算による位相キックバックを用いて，
ユニタリ演算子の固有位相を位相レジスタへ書き込む．
その後，逆量子フーリエ変換を適用することで，位相を計算基底のビット列として読み出す．
有限個の位相量子ビットを用いる場合，測定結果は真の位相に近い値の周囲へ確率的に分布する．

振幅増幅では，良い部分空間に関する反射と初期状態に関する反射を繰り返す．
二つの反射は，良い状態と悪い状態が張る二次元平面上での回転に対応する．
初期成功確率を$a$とすると，単純な反復では$O(1/a)$回を要するのに対し，
振幅増幅では$O(1/\sqrt{a})$回の反復で成功確率を高められる．
また，振幅増幅演算子へ量子位相推定を適用すると，回転角から成功確率を求める振幅推定が得られる．

HHLアルゴリズムでは，量子位相推定によって行列の固有値を固有状態ごとに取り出し，
固有値の逆数に応じた制御回転を行う．
逆量子位相推定によって位相レジスタを初期化した後，
補助量子ビットが成功を示す成分へ振幅増幅を適用できる．
このように，量子位相推定と振幅増幅は，HHLアルゴリズムを構成する独立かつ重要な処理である．

次章では，本章で説明したサブルーチンを前提として，
古典的な線形回帰を線形方程式へ帰着する方法，HHLアルゴリズムによる回帰係数状態の生成，
および量子状態から予測値を評価する方法を扱う．
